
\setlength{\aboverulesep}{-0.7pt}
\setlength{\belowrulesep}{-0pt}
\begin{table*}
	\caption{Comparisons on synthetic dataset (\smvp) evaluated by mean angular error (MAE)~($\downarrow$) with degree and Chamfer distance (CD)~($\downarrow$) in millimeter (mm), respectively. 
	Best and second results in SDF-based and 3DGS-based methods (except PolGS~\citep{han2025polgs}) are highlighted as \colorbox[rgb]{1,0.7,0.7}{1st} and \colorbox[rgb]{1,1,0.4}{2nd}. The time consumed is shown on the far right side of the table.
	}
	\label{table:shape_quantitative_smvp}
	\vspace{-0.5em}
	\centering
	\resizebox{\textwidth}{!}{
\begin{tabular}{lcccccccc|cc|c}
	\toprule
	 & \multicolumn{2}{c}{\sc Hedgehog } & \multicolumn{2}{c}{\sc Squirrel} & \multicolumn{2}{c}{\sc Snail} & \multicolumn{2}{c|}{\sc David} & \multicolumn{2}{c|}{Mean}& \\
	\cmidrule{2-11}
		\multicolumn{1}{c}{\multirow{-2}{*}{Method}} & MAE & CD & MAE & CD & MAE & CD & MAE & CD & MAE & CD& \multicolumn{1}{c}{\multirow{-2}{*}{time}}\\
	\midrule
		\pandora&9.41&9.50  &10.85& 5.88&8.08&10.97&14.75& 4.88&10.77&7.81&10 h\\
	 
	 \nersp& 8.94&6.57 &8.23&  \colorbox[rgb]{1,1,0.4}{3.02}&5.56& \colorbox[rgb]{1,1,0.4}{3.72} &15.38&4.18 &9.53&\colorbox[rgb]{1,1,0.4}{4.37}&10 h\\
	
	 \mvas&\colorbox[rgb]{1,1,0.4}{4.30}&\colorbox[rgb]{1,1,0.4}{4.22} & \colorbox[rgb]{1,1,0.4}{6.10}& 3.73 & \colorbox[rgb]{1,1,0.4}{3.30}& 7.87 & \colorbox[rgb]{1,1,0.4}{8.47}& \colorbox[rgb]{1,1,0.4}{3.21}&\colorbox[rgb]{1,1,0.4}{5.54}&4.76&11 h\\
	 	\gsror &9.72& 4.45&9.25& 4.20&7.24& 5.16&14.02&4.14&10.06&4.49& 1.5 h \\
    
	 \nero &\colorbox[rgb]{1,0.7,0.7}{3.40}&\colorbox[rgb]{1,0.7,0.7}{3.69} & \colorbox[rgb]{1,0.7,0.7}{3.55}& \colorbox[rgb]{1,0.7,0.7}{1.86} & \colorbox[rgb]{1,0.7,0.7}{2.67}& \colorbox[rgb]{1,0.7,0.7}{3.71} &\colorbox[rgb]{1,0.7,0.7}{7.64}& \colorbox[rgb]{1,0.7,0.7}{2.88} &\colorbox[rgb]{1,0.7,0.7}{4.32}&\colorbox[rgb]{1,0.7,0.7}{3.04}&8 h\\
	 
   
    \midrule
    
    \dr &\colorbox[rgb]{1,1,0.4}{12.28}& 12.66&17.18& 11.20 &11.42& 20.70&20.56&7.91&15.36&13.12& 12 min \\

    \relight &14.40& 19.84&\colorbox[rgb]{1,1,0.4}{15.41}&18.97 &\colorbox[rgb]{1,1,0.4}{9.08}& 19.04&\colorbox[rgb]{1,1,0.4}{15.47}& 13.47&\colorbox[rgb]{1,1,0.4}{13.59}&17.83& 25 min \\
	
	\gs &16.50&\colorbox[rgb]{1,1,0.4}{8.82}&21.65& \colorbox[rgb]{1,1,0.4}{9.53}&19.05& \colorbox[rgb]{1,1,0.4}{14.04}&21.56& \colorbox[rgb]{1,1,0.4}{7.39}&19.69&\colorbox[rgb]{1,1,0.4}{9.95}&5 min  \\
    
\polgs~(Ours) &\colorbox[rgb]{1,0.7,0.7}{8.33}&\colorbox[rgb]{1,0.7,0.7}{6.99}&\colorbox[rgb]{1,0.7,0.7}{9.18}& \colorbox[rgb]{1,0.7,0.7}{6.08}&\colorbox[rgb]{1,0.7,0.7}{7.25}&\colorbox[rgb]{1,0.7,0.7}{9.20}&\colorbox[rgb]{1,0.7,0.7}{13.36}& \colorbox[rgb]{1,0.7,0.7}{5.12}&\colorbox[rgb]{1,0.7,0.7}{9.53}&\colorbox[rgb]{1,0.7,0.7}{6.85}&10 min\\				
\midrule
PolGS~\citep{han2025polgs} &10.83& 7.62&11.42& 6.28 &9.64& 10.85
&13.99&5.30&11.47&7.51& 7 min \\

    \bottomrule
\end{tabular}	
}
\end{table*}


\begin{table*}
	\caption{Comparisons on real-world datasets (\rmvp and \pisr) evaluated by mean angular error (MAE)~($\downarrow$) with degree and Chamfer distance (CD)~($\downarrow$) in millimeter (mm), respectively. 
	Best and second results in SDF-based and 3DGS-based methods (except PolGS~\citep{han2025polgs}) are highlighted as \colorbox[rgb]{1,0.7,0.7}{1st} and \colorbox[rgb]{1,1,0.4}{2nd}. The time consumed is shown on the far right side of the table.
	}
	\label{table:shape_quantitative_rmvp_pisr}
	\vspace{-0.5em}
	\centering
	\normalsize
	\resizebox{\textwidth}{!}{
\begin{tabular}{lcccc|cccc|cc|c}
	\toprule
	& \multicolumn{4}{c|}{\rmvp}& \multicolumn{4}{c|}{\pisr}&& \\
	\cmidrule{2-9}
	 & \multicolumn{2}{c}{\sc Shisa} & \multicolumn{2}{c|}{\sc Frog} & \multicolumn{2}{c}{\sc L-Rabbit} & \multicolumn{2}{c|}{\sc S-Rabbit}&\multicolumn{2}{c|}{\multirow{-2}{*}{Mean}}& \\
	\cmidrule{2-11}
		\multicolumn{1}{c}{\multirow{-3}{*}{Method}} & MAE & CD & MAE & CD & MAE & CD & MAE & CD& MAE & CD& \multicolumn{1}{c}{\multirow{-3}{*}{time}}\\
	\midrule
		\pandora&12.93&11.29&15.86&7.88&13.47&16.53&\colorbox[rgb]{1,1,0.4}{11.45}&15.39&13.43&12.77&10 h\\
	 
	 \nersp& 10.79& 7.39& 15.62& 6.68& 13.33& \colorbox[rgb]{1,1,0.4}{7.43}& 15.75&  10.27 &13.87&\colorbox[rgb]{1,1,0.4}{7.94}&10 h\\
	
	 \mvas&8.56& 9.28& 17.63& 7.00&\colorbox[rgb]{1,1,0.4}{11.86}& 13.82&12.81& 12.08&12.72&10.55&11 h\\
	
	\pisr &\colorbox[rgb]{1,0.7,0.7}{8.23}& 6.71& 16.79& \colorbox[rgb]{1,1,0.4}{5.89}&\colorbox[rgb]{1,1,0.4}{11.23}& 10.31&12.18& \colorbox[rgb]{1,1,0.4}{8.49}&\colorbox[rgb]{1,1,0.4}{12.11}&9.20& 0.5 h \\
	\gsror&12.50 &  \colorbox[rgb]{1,1,0.4}{5.30}  &\colorbox[rgb]{1,1,0.4}{14.68}&5.97&18.82&16.40&21.55&19.87&16.89&11.89& 1.5 h \\
	 \nero &\colorbox[rgb]{1,1,0.4}{8.41}& \colorbox[rgb]{1,0.7,0.7}{4.88}& \colorbox[rgb]{1,0.7,0.7}{15.29}& \colorbox[rgb]{1,0.7,0.7}{5.39}& \colorbox[rgb]{1,0.7,0.7}{9.40}& \colorbox[rgb]{1,0.7,0.7}{7.81}& \colorbox[rgb]{1,0.7,0.7}{8.34}& \colorbox[rgb]{1,0.7,0.7}{6.30}&\colorbox[rgb]{1,0.7,0.7}{10.36}&\colorbox[rgb]{1,0.7,0.7}{6.10}&8 h\\
	
    
    \midrule
    
    \dr&19.53&15.87&17.08&27.67&31.23&22.03&30.85&18.20 &24.67&20.94& 12 min \\

    \relight &16.69&14.66&17.85&13.35&29.19&15.66&26.39&20.20&22.53&15.97& 25 min \\
	
	\gs &\colorbox[rgb]{1,1,0.4}{12.79}&\colorbox[rgb]{1,1,0.4}{9.09}&\colorbox[rgb]{1,1,0.4}{16.19}&\colorbox[rgb]{1,0.7,0.7}{7.01}&\colorbox[rgb]{1,1,0.4}{24.77}&\colorbox[rgb]{1,1,0.4}{10.63}&\colorbox[rgb]{1,1,0.4}{22.92}& \colorbox[rgb]{1,1,0.4}{10.78}&\colorbox[rgb]{1,1,0.4}{19.17}&\colorbox[rgb]{1,1,0.4}{9.38}&5 min  \\
    
\polgs~(Ours) &\colorbox[rgb]{1,0.7,0.7}{10.78}& \colorbox[rgb]{1,0.7,0.7}{7.16} &\colorbox[rgb]{1,0.7,0.7}{14.63}&\colorbox[rgb]{1,1,0.4}{7.15}&\colorbox[rgb]{1,0.7,0.7}{11.51}&\colorbox[rgb]{1,0.7,0.7}{7.52}&\colorbox[rgb]{1,0.7,0.7}{12.30}&\colorbox[rgb]{1,0.7,0.7}{8.58} &\colorbox[rgb]{1,0.7,0.7}{12.31}&\colorbox[rgb]{1,0.7,0.7}{7.60}&10 min\\				
\midrule
PolGS~\citep{han2025polgs} &10.88& 7.76& 15.03& 7.48& 11.55& 7.76& 12.73& 9.29& 12.55&8.07& 7 min \\
    \bottomrule
\end{tabular}	
}
\end{table*}
\section{Experiments}
\label{sec:exp}
To evaluate the performance of our method, we conduct the following experiments: 1) Shape reconstruction 
on synthetic dataset, 2) Shape reconstruction on real-world dataset, 3) Ablation study and 4) Radiance decomposition.

% \begin{figure}
% 	\centering
% 	\includegraphics[width=\linewidth]{image/cameras_position.pdf}
% 	\caption{Visualization of view distributions of different datasets. (a) Input views display on \smvp. (b) Input views display on \rmvp. (c) Input views display on \pandora. (d) Input views display on \pisr.}
% 	\label{fig:camera}
% \end{figure}




% \begin{figure*}
% 	\LARGE
% 	\includegraphics[width=\linewidth]{image/synthetic_2.pdf}
	
% 	\caption{Qualitative comparisons on surface normal in \smvp.
% %		 The mean angular error map has shown in range of $0^\circ$ to $20^\circ$.
% 		%		 More results can be found in the supplementary material.
% 	}
% 	\vspace{-0em}
% 	\label{fig:syn_normal_2}
% \end{figure*}

\vspace{-1.4em}\paragraph{Dataset\quad}  We use four datasets to evaluate our method: the synthetic dataset \smvp, the real-world dataset \rmvp,
\pandora and \pisr, where \pandora can only be used for qualitative evaluation due to the lack of ground truth.

% (1) \smvp. The input images are all at a resolution of $512 \times 512$ pixels. We select 36 different views around each object to ensure adequate coverage. The specific input views are displayed in \fref{fig:camera}~(a).

% (2) \rmvp. The original input resolution is $1024 \times 1224$, which we resize to $512 \times 612$ during training and evaluation. We use 31 views to train on each object. The specific input views are displayed in \fref{fig:camera}~(b).

% (3) \pandora. The original input resolution is $2048 \times 2448$, which we resize to $512 \times 612$ during training and evaluation. We use all 35 views to train on each object. The specific input views are displayed in \fref{fig:camera}~(c).

% (4) \pisr. The original input resolution is $2048 \times 2448$, which we resize to $512 \times 612$ during training and evaluation. We use 20 views for training and evaluation. The specific input views are displayed in \fref{fig:camera}~(d).
\vspace{-1.0em}
\paragraph{Baselines\quad} 
We compare our method with state-of-the-art techniques, including 3DGS methods \gs, \dr, and \relight, and SDF-based methods \nero, \mvas, \pandora, \nersp, \pisr, and \gsror(which surface strongly depends on SDF representation). 
All of these methods, except \gs, can handle reflective surfaces. \mvas, \pandora, \pisr, and \nersp utilize 
polarimetric information to reconstruct shapes. Specifically, we do not add the normal prior in \gs among the whole experiments for fair comparisons.
\vspace{-1em}
\begin{figure*}
	\LARGE
	\begin{overpic}
    [width=\textwidth]{image/synthetic}
    \put(8.8,94.5){\color{black}{\fontsize{6pt}{1pt}\selectfont {\cite{liu2023nero}}}}
    \put(18.3,94.5){\color{black}{\fontsize{6pt}{1pt}\selectfont {\cite{cao2023multi}}}}
    \put(27,94.5){\color{black}{\fontsize{6pt}{1pt}\selectfont {\cite{dave2022pandora}}}}
    \put(36.5,94.5){\color{black}{\fontsize{6pt}{1pt}\selectfont {\cite{han2024nersp}}}}
    \put(45.5,94.5){\color{black}{\fontsize{6pt}{1pt}\selectfont {\cite{zhu2025gs}}}}
    \put(54.2,94.5){\color{black}{\fontsize{6pt}{1pt}\selectfont {\cite{dai2024high}}}}
    \put(63.7,94.5){\color{black}{\fontsize{6pt}{1pt}\selectfont {\cite{ye20243d}}}}
    \put(72.2,94.5){\color{black}{\fontsize{6pt}{1pt}\selectfont {\cite{gao2023relightable}}}}
    \put(81.5,94.5){\color{black}{\fontsize{6pt}{1pt}\selectfont {\cite{han2025polgs}}}}
	\end{overpic}
	\caption{Qualitative comparisons on surface normal in \smvp, where our method can outperform existing methods based on the same representation and achieves comparable results with SDF-based methods.
%		 The mean angular error map has shown in range of $0^\circ$ to $20^\circ$.
		%		 More results can be found in the supplementary material.
	}
	\vspace{-0.5em}
	\label{fig:syn_normal}
    
\end{figure*}
\paragraph{Evaluation metrics\quad}
We use the Chamfer Distance (CD) to evaluate the shape recovery performance and 
the mean angular error (MAE) to evaluate the quality of surface normal estimations. 3DGS-based methods provide point cloud results and we use Poisson surface reconstruction to generate the mesh for the evaluation.
\vspace{-1.2em}
\begin{figure*}
	\huge 
	\begin{overpic}
    [width=\textwidth]{image/real-world}
    \put(8.8,57.8){\color{black}{\fontsize{5.5pt}{1pt}\selectfont {\cite{liu2023nero}}}}
    \put(16.5,57.8){\color{black}{\fontsize{5.5pt}{1pt}\selectfont {\cite{cao2023multi}}}}
    \put(25,57.8){\color{black}{\fontsize{5.5pt}{1pt}\selectfont {\cite{dave2022pandora}}}}
    
    \put(33.7,57.8){\color{black}{\fontsize{5.5pt}{1pt}\selectfont {\cite{chen2024pisr}}}}
    \put(42,57.8){\color{black}{\fontsize{5.5pt}{1pt}\selectfont {\cite{zhu2025gs}}}}
    \put(50,57.8){\color{black}{\fontsize{5.5pt}{1pt}\selectfont {\cite{han2024nersp}}}}
    
    \put(58,57.8){\color{black}{\fontsize{5.5pt}{1pt}\selectfont {\cite{dai2024high}}}}
    \put(66.5,57.8){\color{black}{\fontsize{5.5pt}{1pt}\selectfont {\cite{ye20243d}}}}
    \put(74.5,57.8){\color{black}{\fontsize{5.5pt}{1pt}\selectfont {\cite{gao2023relightable}}}}
    \put(83.0,57.8){\color{black}{\fontsize{5.5pt}{1pt}\selectfont {\cite{han2025polgs}}}}
	\end{overpic}
	\caption{Qualitative comparison shape on \pandora, \rmvp and \pisr, where PolGS produces similar quality of reconstruction mesh compared to SDF-based methods. }
	\label{fig:PANDORA-result}
    \vspace{-0.5em}
\end{figure*}
\paragraph{Implementation details\quad}
We conduct the experiments on an NVIDIA RTX 4090 GPU with 24GB memory. The training of our model is implemented in 
PyTorch 1.12.1~\citep{paszke2019pytorch} using the Adam optimizer~\cite{kingma2014adam}. The specific learning rates of different components
in Gaussian kernels are set the same as \gs. We adopt a warm-up strategy to train the model, and we add the polarimetric information and deferred rendering after 1000 iterations.


\subsection{Reconstruction results on synthetic dataset}

In \fref{fig:syn_normal} and \Tref{table:shape_quantitative_smvp}, we compare the shape recovery performance of various methods on the \smvp, which contains four objects with spatially varying and reflective properties. However, it is worth noting that SDF-based methods, such as \nero, \gsror, \mvas, \pandora, and \nersp, outperform 3DGS methods in terms of surface representation. This advantage is largely due to their superior ability to model complex surface details. Among the 3DGS methods, \gs struggles with reflective surfaces, while \dr and \relight still has difficulty in representation of surface normal accurately.
In contrast, \polgs effectively leverages polarimetric information, significantly improving geometric surface performance. Due to the inadequate point cloud generation by other 3DGS methods, they fail to produce reasonable mesh results using Poisson surface reconstruction. Our method, however, achieves the lowest mean Chamfer Distance across the synthetic \smvp dataset, reinforcing the trends observed in surface normal estimations and confirming that our approach delivers the best performance among the 3DGS techniques.

\subsection{Reconstruction results on real-world dataset}
We further evaluate the reconstruction performance on the \pandora, \rmvp,  and \pisr datasets, with qualitative results illustrated in \fref{fig:PANDORA-result} and quantitative results detailed in \Tref{table:shape_quantitative_rmvp_pisr}. In real-world scenarios, our method produces reconstructions that align more closely with SDF-based approaches, while significantly outperforming other 3DGS-based methods in terms of reconstruction quality. For example, in the {\sc Vase} case, our approach accurately estimates the shape of a ceramic surface in just $10$ minutes, achieving results closer to those generated by SDF-based methods. 
Additionally, the {\sc Frog} sample highlights our method's ability to reconstruct objects with rough glossy surfaces, showing the robustness and generalization capability of our method. 
Futhermore, texure-less object in \pisr dataset also demonstrates the advantage of our method in handling objects with limited RGB cues, where the polarimetric information provides critical constraints for accurate shape recovery.
These results collectively demonstrate the effectiveness of our approach in handling diverse real-world objects with varying surface properties. 




\begin{figure*}
	\centering
	\includegraphics[width=\textwidth]{image/ablation_table.pdf}
	\caption{ Quantitative ablation results on \smvp, where we evaluate the effect of polarimetric supervision under different input views. 
	% The results show that polarimetric information provides more significant benefits when the number of views is limited, and the combination of photometric polarization loss and tangent-space consistency loss yields the best performance across all settings.
	}
	\label{fig:ablation_table}
	
\end{figure*}
\begin{figure}
	\centering
	\includegraphics[width=\linewidth]{image/ablation.pdf}
	\caption{
Qualitative ablation results on \pisr, where we compare the reconstruction quality with and without polarimetric supervision. 
% The results demonstrate that polarimetric cues significantly enhance the reconstruction quality, by providing additional constraints that help resolve local shape ambiguities.
	}
	\label{fig:ablation}
	\vspace{-1em}
\end{figure}
\subsection{Ablation study}\label{sec:ablation}

In this section, we conduct two ablation studies to validate key design choices in our method.

\paragraph{Effect of polarimetric information\quad}
We first study how polarimetric cues contribute to surface reconstruction. Our polarimetric supervision consists of two terms: a photometric polarization loss $\mathcal{L}_{pol}$ and a tangent-space consistency loss $\mathcal{L}_{tsc}$. We perform controlled experiments on \smvp by varying both the number of input views and the enabled loss terms (w/o polarimetric loss, w/ $\mathcal{L}_{pol}$ only, and w/ both).

\Fref{fig:ablation_table} shows a clear trend: when view coverage is dense, RGB photometric cues already provide strong constraints, and the additional gain from $\mathcal{L}_{pol}$ becomes marginal. In contrast, $\mathcal{L}_{tsc}$ yields a more consistent improvement, indicating that enforcing stable polarization-derived constraints is particularly beneficial even in well-observed settings. As the number of views decreases, the advantage of polarimetric information becomes increasingly evident, and combining $\mathcal{L}_{pol}$ with $\mathcal{L}_{tsc}$ leads to the most robust reconstructions under sparse-view conditions.

We further provide a qualitative evaluation on \pisr to highlight the practical value of polarization in \fref{fig:ablation}. While the RGB appearance lacks discernible texture, the polarization channels $s_1$ and $s_2$ and AoP information exhibit meaningful variations that provide additional geometric cues for resolving shape ambiguities. The results show that both concave and convex surface distortions are more accurately reconstructed when polarization information is employed.

\begin{figure}
	\centering
	\includegraphics[width=\linewidth]{image/visibility_ablation.pdf}
	\caption{Ablation results on the visibility mask threshold.}
	\label{fig:visibility_ablation}
	\vspace{0em}
\end{figure}

\paragraph{Visibility mask threshold\quad}
\label{exp:visibility}
We then validate the threshold $\tau$ used to binarize the depth-guided visibility mask in \Eref{eq:visibility}. We sweep different thresholds and compare the resulting masks to the SDF-based visibility mask under ground-truth normal supervision. 

As shown in \fref{fig:visibility_ablation}, directly using the SDF-based visibility masks yields a surface normal MAE of 10.594 degrees in \polgs. With depth-guided masks, performance varies with $\tau$, and $\tau = 0.010$ provides the best trade-off: an accuracy of 0.9489 against the GT visibility mask and an MAE of 10.784 degrees. Without any visibility mask, the MAE degrades to 11.141 degrees, confirming the benefit of incorporating visibility information. We observe the same trend in other datasets and therefore use $\tau = 0.010$ for all experiments in the paper.






%
\subsection{Comparison of radiance decomposition}
\Fref{fig:env_compare} presents a comparison of diffuse and specular component decomposition generated by \pandora, \dr, and \polgs. Here, \pandora utilizes the IDE~\cite{verbin2022ref} structure to produce environment map results, whereas \polgs adopts \dr's methods by employing a Cubemap encoder for the same purpose. 
Compared with \dr, \polgs leverages additional polarimetric information to effectively constrain and disambiguate diffuse and specular components. 
Notably, our results closely align with \pandora's, demonstrating improved radiance decomposition ability.
\begin{figure}
	\begin{overpic}[width=\linewidth, trim={0pt 0pt 0pt 0pt}]{image/env.pdf}
		
		% \put(75,43){\color{black}{\fontsize{8pt}{1pt}\selectfont {\polgs}}}
		% \put(10,43){\color{black}{\fontsize{8pt}{1pt}\selectfont {PANDORA}}}
		% \put(41.5,43){\color{black}{\fontsize{8pt}{1pt}\selectfont {3DGS-DR}}}
	\end{overpic}
	\vspace{-1.5em}
	\caption{Separation of diffuse and specular components with \pandora, \dr and \polgs.}
	\label{fig:env_compare}
	\vspace{-0em}
\end{figure}
